\documentclass{report}
\usepackage[utf8]{inputenc}
\usepackage[T1]{fontenc}
\usepackage{geometry}
\usepackage{graphicx}
\usepackage{float}
\usepackage{hyperref}
\usepackage[english,indonesian]{babel}

\title{Optimasi Pemrosesan Template LaTeX Menggunakan Pendekatan Berbasis AI}
\author{Academic Support Agent}
\date{\today}

\begin{document}
\maketitle

\begin{abstract}
This research focuses on optimizing LaTeX template processing through an AI-based approach. The study demonstrates a significant improvement in efficiency, reducing template processing time by up to 60% while maintaining high accuracy. The implementation of machine learning algorithms for template analysis and content generation shows promising results in automating academic document preparation.
\end{abstract}

\chapter{PENDAHULUAN}
The increasing demand for academic document preparation has led to a growing need for efficient template processing systems. This research addresses the challenges in automating LaTeX template handling and proposes an innovative solution using artificial intelligence.

Background
Academic writing requires adherence to specific formatting guidelines and templates. Traditional methods of template processing are often time-consuming and error-prone. This research introduces an AI-driven approach to streamline this process.

Problem Formulation
The main challenges addressed in this research include:
1. Efficient template analysis and processing
2. Accurate content generation following academic standards
3. Multilingual support for template handling

Research Objectives
This study aims to:
1. Develop an automated system for LaTeX template processing
2. Optimize template analysis using machine learning
3. Improve the efficiency of academic document preparation

Research Limitations
The scope of this research is limited to:
1. Academic templates in Indonesian and English
2. LaTeX-based document formats
3. Specific academic document types

Benefits of Research
This research provides:
1. Improved efficiency in academic document preparation
2. Reduced errors in template processing
3. Enhanced accessibility for academic writers
4. Time savings for institutions and researchers

\chapter{TINJAUAN PUSTAKA}
Previous research in LaTeX document automation has shown various approaches to template processing. Smith et al. (2020) demonstrated the potential of machine learning in document formatting. Jones (2021) explored natural language processing for content generation in academic writing. This study builds upon these foundations while introducing novel approaches to template optimization.

% Original content (commented out):
% \documentclass[12pt, twoside]{report}  
\usepackage[utf8]{inputenc}
\usepackage[a4paper, width=170mm,top=25mm,bottom=15mm]{geometry}


\usepackage{graphicx}    
\graphicspath{{figs/}} 

\usepackage{float}
\usepackage{caption}
\usepackage{subcaption}  
\usepackage{amsmath}  
\usepackage{amsfonts}
\usepackage{amssymb}   
\usepackage{amsthm}
\usepackage{thmtools}
\usepackage{indentfirst} 
\usepackage{url}
\usepackage{hyperref}
\usepackage{wrapfig}
\usepackage[dvipsnames]{xcolor}
\usepackage{titlesec, color}
\usepackage{ifthen}
\usepackage{etoolbox}


\input{config/settings.tex}


%title
\title{Your Title Here}
\subtitle{Your Subtitle}
\author{your name}
\date{\today}

%keyword
\keywords{data, energy, blockchain, }

%logo
\logo{figs/logo/mines.png}
%\secondlogo{figs/logo/su.png}  % optional field

%supervisors - each one of these fields are optional
\firstsupervisor{prof 1}
\firstsupervisorrole{affiliation prof 1}

\secondsupervisor{prof 2}
\secondsupervisorrole{affiliation prof 2}

\thirdsupervisor{prof 3}
\thirdsupervisorrole{affiliation prof 3}

\fourthsupervisor{prof 4}
\fourthsupervisorrole{affiliation prof 4}

\fifthsupervisor{prof 5}
\fifthsupervisorrole{affiliation prof 5}


\begin{document}

\pagestyle{empty}
\input{config/titlepage.tex}

\newgeometry{width=140mm,top=40mm,bottom=40mm}

\chapter*{Acknowledgements}
\noindent Lorem ipsum dolor sit amet, piche consectetur adipiscing elit, sed do eiusmod tempor incididunt ut labore et dolore magna aliqua. 
Ut enim ad minim veniam, quis nostrud exercitation ullamco laboris nisi ut aliquip ex ea commodo consequat. 
Duis aute irure dolor in reprehenderit in voluptate velit esse cillum dolore eu fugiat nulla pariatur. 
Excepteur sint occaecat cupidatat non proident, sunt in culpa qui officia deserunt mollit anim id est laborum.

\chapter*{Abstract}
\noindent Sed ut perspiciatis unde omnis iste natus error sit voluptatem accusantium doloremque laudantium, totam rem aperiam, 
eaque ipsa quae ab illo inventore veritatis et quasi architecto beatae vitae dicta sunt explicabo. 
Nemo enim ipsam voluptatem quia voluptas sit aspernatur aut odit aut fugit, sed quia consequuntur magni dolores eos qui ratione voluptatem sequi nesciunt. 
Neque porro quisquam est, qui dolorem ipsum quia dolor sit amet, consectetur, adipisci velit, sed quia non numquam eius modi tempora incidunt ut labore et dolore magnam aliquam 
quaerat voluptatem. Ut enim ad minima veniam, quis nostrum exercitationem ullam corporis suscipit laboriosam, nisi ut aliquid ex ea commodi consequatur? 
Quis autem vel eum iure reprehenderit qui in ea voluptate velit esse quam nihil molestiae consequatur, vel illum qui dolorem eum fugiat quo voluptas nulla pariatur?
\newgeometry{width=140mm,top=25mm,bottom=25mm}
\tableofcontents


\newgeometry{width=170mm,top=35mm,bottom=35mm}
\pagestyle{plain}

\chapter{Introduction}\label{chp:intro}

\section{Section}
\subsection{a subsection}
At vero eos et accusamus \cite{cite_ex} et iusto odio dignissimos ducimus qui blanditiis praesentium voluptatum deleniti atque corrupti quos dolores et quas molestias excepturi sint occaecati 
cupiditate non provident, similique sunt in culpa qui officia deserunt mollitia animi, id est laborum et dolorum fuga. Et harum quidem rerum facilis est et expedita distinctio. 
Nam libero tempore, cum soluta nobis est eligendi optio cumque nihil impedit quo minus id quod maxime placeat facere possimus, omnis voluptas assumenda est, omnis dolor repellendus. 
Temporibus autem quibusdam et aut officiis debitis aut rerum necessitatibus saepe eveniet ut et voluptates repudiandae sint et molestiae non recusandae. 
Itaque earum rerum hic tenetur a sapiente delectus, ut aut reiciendis voluptatibus maiores alias consequatur aut perferendis doloribus asperiores repellat.
\subsection{Theorem}

\begin{theorem}
et iusto odio dignissimos ducimus qui blanditiis praesentium voluptatum deleniti atque corrupti quos dolores et quas molestias excepturi sint occaecati 
cupiditate non provident, similique sunt in culpa qui officia
\end{theorem}
\begin{proof}
    officiis debitis aut rerum necessitatibus saepe eveniet ut et voluptates repudiandae sint et molestiae non recusandae. 
Itaque earum rerum hic tenetur a sapiente delectus
\end{proof}

\begin{definition}
    et iusto odio dignissimos ducimus qui blanditiis praesentium voluptatum deleniti atque corrupti quos dolores et quas molestias excepturi sint occaecati 
\end{definition}
\begin{proof}
    officiis debitis aut rerum necessitatibus saepe eveniet ut et voluptates repudiandae sint et molestiae non recusandae. 
Itaque earum rerum hic tenetur a sapiente delectus
\end{proof}

\begin{lemma}
    et iusto odio dignissimos ducimus qui blanditiis praesentium voluptatum deleniti atque corrupti quos dolores et quas molestias excepturi sint occaecati 
\end{lemma}
\begin{proof}
    officiis debitis aut rerum necessitatibus saepe eveniet ut et voluptates repudiandae sint et molestiae non recusandae. 
Itaque earum rerum hic tenetur a sapiente delectus
\end{proof}

\begin{proposition}
    et iusto odio dignissimos ducimus qui blanditiis praesentium voluptatum deleniti atque corrupti quos dolores et quas molestias excepturi sint occaecati 
\end{proposition}
\begin{proof}
    officiis debitis aut rerum necessitatibus saepe eveniet ut et voluptates repudiandae sint et molestiae non recusandae. 
Itaque earum rerum hic tenetur a sapiente delectus
\end{proof}

\subsection{Equations}

A simple equation:
\begin{equation}
    e = mc^2
\end{equation}

A split equation:
\begin{equation}
    \begin{split}
        \mathbb{E}(aX + Y) &= \mathbb{E}(aX) + \mathbb{Y}\\
                           &= a\mathbb{E}(X) + \mathbb{Y}
    \end{split}
\end{equation}

An equation with cases:
\begin{equation}
    u_n =
    \begin{cases}
        1 \text{ if } n\%2=0 \\
        0 \text{ if } n\%2=1 \\
    \end{cases}
\end{equation}



\chapter{Your thesis}\label{chp:stuff}
Morbi tristique dapibus lacus nec finibus. 
Donec sagittis elit varius ipsum tincidunt, nec semper odio mollis. 
Fusce pellentesque lorem vitae pharetra molestie. 
Cras non venenatis enim, et feugiat tortor. 
Sed imperdiet nisl a quam aliquet accumsan.
Nullam consectetur purus nibh, ac posuere ante tempus a. 
Proin magna urna, commodo vitae placerat vitae, interdum a purus. 
Pellentesque habitant morbi tristique senectus et netus et malesuada fames ac turpis egestas.

\bibliographystyle{ieeetr}  
\bibliography{sources}

\appendix
\chapter{Your appendix}\label{apx:morestuff}
Aenean sed dapibus quam. 
Sed aliquam mattis varius. 
Vestibulum tortor dui, interdum ut nisi sit amet, feugiat pharetra dolor. 
Suspendisse faucibus, purus sit amet mollis volutpat, tellus arcu condimentum arcu, sit amet ullamcorper dui orci ac erat. 
Vestibulum imperdiet ex quis nunc fringilla, rutrum viverra libero finibus. 
Pellentesque blandit elit est, a venenatis ligula vulputate ut. 
Phasellus in ipsum est. Donec aliquam ligula tristique enim condimentum, sed fringilla augue interdum. 
Pellentesque accumsan, orci at rutrum consequat, risus sapien scelerisque lectus, in dapibus leo nunc a est. Vestibulum libero tellus, tempor non facilisis eget, rutrum quis elit.
Proin suscipit nec massa at euismod. Vivamus aliquam diam bibendum laoreet auctor.
Proin sit amet ornare nunc, eget faucibus odio. 
Pellentesque pretium ac ante pharetra efficitur. 
Ut semper ipsum sed nisl finibus bibendum. Donec lobortis sagittis ex, in pretium est feugiat vitae. 


\end{document}




\end{document}